\section{Stopping, Withdrawing, and Changing Your Mind}

\subsection{Four routes out, and what each does to your record}

\draftinstr{Roughly 1,200 characters introducing Figure 22. Participants are
routinely told they may withdraw at any time and are almost never told what
withdrawal does to the data already collected. The honest answer depends on
which route out is taken, and there are four: stop the drug only; stop study
procedures but keep follow-up; full withdrawal; and withdrawal of consent for
data use as well. Adapt from \texttt{inputs/\pfb phase-\pfb 1-\pfb trial-\pfb protocol.zip},
\texttt{sections/\pfb sec-\pfb 07-\pfb discontinuation.tex}, and cite \texttt{ecfr2026part50}
and \texttt{NCIConsent24}.}

\begin{pafig}[plantuml-type: activity with swimlanes, @startuml]
\node[umlkey] (a) at (0,0) {Figure 22 slot\\withdrawal activity};
\node[umlbox] (b) at (0,-1.9) {\ttfamily\tiny Source\\ plantuml/\\ fig-\pfb 22-\pfb withdrawal-\pfb activity.puml};
\draw[umlarrow] (a) -- (b);
\end{pafig}
\vspace{-0.7cm}
\figcaption{Figure 22. The four routes out of the study drawn as a swimlaned activity, with
one lane for the participant, one for the site team, and one for the data and specimens,\\
so each route's consequence for the record is visible beside the route itself.}

\draftinstr{Replace the Figure 22 slot with the PlantUML-type activity from
\texttt{plantuml/\pfb fig-\pfb 22-\pfb withdrawal-\pfb activity.puml}. Three vertical
\texttt{umlpkg} swimlanes at $x = -5.0, 0.6, 6.6$ sharing one top and one bottom
edge; actions as \texttt{umlstate} boxes at 2.2 cm pitch; one \texttt{mmdec}
decision at $(0.6,-3.0)$; fork and join as \texttt{umlbar} rectangles; four
\texttt{umlfinal} terminals on one rank; cross-lane edges with right-angle
routing, the two that change rank at looseness 0.85.}

\vspace{0.30em}

\noindent\begin{tabularx}{\textwidth}{L{2.1cm} L{4.0cm} Y}
\toprule
\textbf{Route} & \textbf{What stops} & \textbf{What happens to data and specimens} \\
\midrule
A & Daraxonrasib only & Everything already collected stays in the safety analysis; new data continues to be collected \\
B & Study procedures, follow-up continues & Existing data retained; new data limited to survival status; specimens retained unless you ask for destruction \\
C & All study participation & Data already collected cannot be removed from the safety analysis; nothing new is collected; specimens destroyed on request \\
D & Participation and future data use & Identifiable data segregated from new analysis; aggregate safety data already filed with the FDA cannot be recalled \\
\bottomrule
\end{tabularx}

\vspace{0.30em}

\draftinstr{Do not soften routes C and D. Both statements are true and
unwelcome, and omitting them would make routes A and B less believable.}

\subsection{Consent as a state, not a signature}

\draftinstr{Roughly 1,100 characters introducing Figure 23. A consent form
implies a one-time act. In a Physical AI trial the thing consented to can change
while the participant is still enrolled, because an AI-enabled device can be
updated. This protocol freezes the software version at consent and forces a
re-consent on any change, so the participant is never operated on by software
they have not been told about. Cite \texttt{FDATrans2024} and
\texttt{FDARealW2025}.}

\begin{pafig}[mermaid-type: stateDiagram-v2, composite]
\node[mmgoal] (a) at (0,0) {Figure 23 slot\\consent lifecycle};
\node[mmin] (b) at (0,-2.0) {\ttfamily\tiny Source\\ mermaid/\\ fig-\pfb 23-\pfb consent-\pfb lifecycle.md};
\draw[mmedge] (a) -- (b);
\end{pafig}
\vspace{-0.7cm}
\figcaption{Figure 23. Consent drawn as a state machine with a nested enrolled composite,
withdrawal reachable from every state, and a re-consent transition the participant did\\
not cause, triggered by a software version change, a new risk, or a protocol amendment.}

\draftinstr{Replace the Figure 23 slot with the mermaid-type state diagram from
\texttt{mermaid/\pfb fig-\pfb 23-\pfb consent-\pfb lifecycle.md}. \texttt{umlinit} at $(0,1.4)$,
two \texttt{umlfinal} discs at $y=-13.0$; the enrolled composite as a
\texttt{umlpkg} fitted over five \texttt{umlstate} substates at 1.4 cm pitch;
transitions out of the composite leave the composite border rather than a
substate, at looseness 0.9; the two notes as \texttt{umlnote} boxes with 0.5pt
dashed leaders.}

\subsection{Discontinuation the study may initiate}

\draftinstr{Roughly 900 characters. Distinguish clearly between the participant
stopping and the study stopping the participant. Populate the triggers from
\texttt{inputs/\pfb phase-\pfb 1-\pfb trial-\pfb protocol.zip},
\texttt{sections/\pfb sec-\pfb 07-\pfb discontinuation.tex} \S7.2, and state that in every case
the participant is told the reason, in writing, and returns to standard care.}

\vspace{0.30em}

\noindent\begin{tabularx}{\textwidth}{L{5.0cm} Y}
\toprule
\textbf{Trigger} & \textbf{What follows for you} \\
\midrule
Dose-limiting toxicity at your level & Drug held or discontinued; surgical follow-up and every safety assessment continue \\
Disease progression before surgery & Curative-intent resection is no longer the right operation; you are referred back for systemic therapy \\
Intercurrent illness precluding surgery & Participation ends; standard-of-care pathway resumes without penalty \\
Study-wide pause or halt & You are told, in writing, within the timeframe the protocol specifies; follow-up continues \\
Withdrawal of the IND or the IDE & Participation ends; the sponsor's obligation for research-related injury care continues \\
\bottomrule
\end{tabularx}

\vspace{0.30em}

\subsection{Lost to follow-up}

\draftinstr{Roughly 600 characters. Define what counts as lost to follow-up,
state the contact attempts the protocol requires before that determination, and
state that a participant who reappears is resumed rather than excluded. Source
from \S7.3 of the parent protocol.}
